\subsection*{期值与等待}

期值命名传输结果。异步传输的期值还可以传给 \co@wait@，用来等待完成。

\co@wait@ 可以同时等待多个异步期值；同步传输已经完成，不能再用其期值等待。

\begin{lstlisting}[language=croqtile,caption={同步与异步期值的等待}]
f0 = dma.copy.async input0 => shared;
f1 = dma.copy.async input1 => shared;
f2 = dma.copy input2 => shared;
...
wait f0, f1;  // Multiple wait
wait f2;      // Error: cannot wait on a sync-dma
\end{lstlisting}

同步传输返回期值，是为了让后续代码取得目的缓冲区。主要成员如下：

\begin{itemize}[leftmargin=1.6em]
\item \co@.span@ 取出 DMA 目的缓冲区的 \co@mdspan@。
\item \co@.data@ 取出 DMA 目的缓冲区的引用（带形状数据）。
\end{itemize}

下面通过期值访问目的缓冲区：

\begin{lstlisting}[language=croqtile,caption={通过期值访问数据与形状}]
f0 = dma.copy.async input0 => shared;
f1 = dma.copy input2 => shared;
...
local f32 [f0.span] buffer;
call device_kernel(f1.data, |f1.span|);
\end{lstlisting}

\co@.span@ 返回目的形状，\co@.data@ 返回目的数据。示例中的元素计数运算还可以取得目的区域的总元素数。
